\clearpage
\fancyhead[L]{\small PALLAS.}
\hypertarget{native-267-lower}{}
\begin{center}
PALLAS.

\rule{0.5\linewidth}{0.4pt}

Monatliche Correspondenz. Vol. XI. pp. 376–383. April 1805.

\rule{0.5\linewidth}{0.4pt}
\end{center}

It has been impossible to observe \emph{Pallas} for several months now. Thanks to Dr. \textsc{Gauss}’s diligence, however, we can give our readers some further theoretical information about this new planet. It will be especially useful in finding the planet again when it next becomes visible, which is probably still some time away. Only a lack of material to work from can set limits to Dr. \textsc{Gauss}’s tireless industry and sustained efforts to perfect the calculated orbits of the planets. He expressed his dissatisfaction more than once at the scarcity of observations of Pallas, which made it impossible for him to undertake a further correction of the seventh set of elements of this planet. Admittedly, observing this faint celestial body presented various difficulties, and the observations made by \textsc{Olbers}, \textsc{Oriani}, \textsc{David} and at our observatory are the only ones made during the whole of the past year. The opposition of Pallas observed at \emph{Brera} on 30 August 1804 was therefore very welcome to him. At that time the mean error of the seventh set of elements was $-7'\,28''3$ in longitude and $+2'\,14''7$ in latitude. He based the following eighth set of elements of Pallas on this and on all available earlier observations.
\SourceLeaf{268}
\begin{center}
\begingroup\fontsize{9}{11}\selectfont\setlength{\tabcolsep}{3pt}\renewcommand{\arraystretch}{1.1}
\begin{tabular}{@{}l|l@{}}
\begin{tabular}[t]{@{}lr@{}}
\shortstack[l]{Epoch, Seeberg\\meridian}  & $\left\{\begin{array}{lr}1803\ .&221^\circ\ 31'\ 23''2\\1804\ .&299\ 58\ 38.1\end{array}\right.$ \\
\shortstack[l]{Daily tropical\\motion . . .} & $771''6802$ \\
\shortstack[l]{Aphelion 1803\\ . . . \raisebox{.3ex}{.} . . .} & $301^\circ\ 1'\ 44''1$ \\
\end{tabular} &
\begin{tabular}[t]{@{}lr@{}}
\shortstack[l]{Ascending node\\1803 . . . .} & $172^\circ\ 29'\ 6''8$ \\
\shortstack[l]{Eccentricity\\ . . . . . . . . .} & $0.246101$ \\
\shortstack[l]{Logarithm of the\\semiaxis . . .} & $0.4417647$ \\
\shortstack[l]{Inclination of\\the orbit . . . . . . .} & $34^\circ\ 37'\ 43''2$ \\
\end{tabular} \\
\end{tabular}\endgroup
\end{center}

\indent From these elements Dr. \textsc{Gauss} also calculated the following ephemeris for the geocentric path of \emph{Pallas} [from 28 July 1805 to 30 April 1806, giving the positions in the same way as for Ceres in the same years, \hyperlink{native-267}{p. 262} of this volume.]

\Needspace{19\baselineskip}
—\quad—\quad—

Dr. \textsc{Gauss} wrote to us that this ephemeris was again based on purely elliptical elements. He thought it was still inadvisable to take perturbations into account for Pallas. For one thing, a pure ellipse still represents all observations to date quite well. For another, the short period during which Pallas was visible had not allowed its orbit to be determined accurately enough to expect that working out all the perturbation equations would yield a degree of accuracy that would repay the labor involved. —\quad—\quad—\quad—\quad—\quad— If the mass of Ceres were somewhat greater, it might cause substantial perturbations in the orbit of Pallas, partly because the two planets are so close to one another and partly because their mean motions are nearly equal.

—\quad—\quad—

Inspector \textsc{Harding}, to whom Dr. \textsc{Gauss} immediately sent the new ephemeris for Pallas given above, was engaged in the very worthwhile work of drawing up a star chart for its path. Unfortunately, he will often lack measured star positions here, since in these southern zones the Histoire céleste gives him less help than this valuable collection ordinarily provides.

\hfill[\textsc{von Zach.}]
